% CLAUDE REPORT::WR-BRIDGE-0002::2026-11-14T11:20:00-08:00
\documentclass[11pt]{article}
\usepackage[margin=1in]{geometry}
\usepackage{amsmath,amssymb,amsthm,booktabs}
\usepackage{xurl}
\usepackage[hidelinks]{hyperref}
\newtheorem{theorem}{Theorem}
\newtheorem{proposition}[theorem]{Proposition}
\newtheorem{lemma}[theorem]{Lemma}
\newcommand{\SAT}{\mathrm{SAT}}
\newcommand{\Pc}{\mathsf P}
\newcommand{\NP}{\mathsf{NP}}
\newcommand{\poly}{\operatorname{poly}}
\newcommand{\WR}{\overline W_R}
\newcommand{\Down}{\mathrm{Down}}
\newcommand{\Base}{\mathrm{Base}}
\title{Attempting Both Bridge Targets:\\
An Inductive-Invariant Form of the EF Bridge, and Why Weak Systems Fail}
\author{Claude routine (worst-case derivation directive), run WR-BRIDGE-0002}
\date{14 November 2026, 11:20 America/Los\_Angeles}
\begin{document}
\maketitle

\section*{Status}
\textbf{P versus NP remains OPEN.} This run attempts the two targets left by
WR-BRIDGE-0001:
(A) $\SAT\in\Pc\Rightarrow$ some polynomial-time SAT decider has short
Extended Frege (EF) proofs of its correctness;
(B) a system $R$ strictly between Resolution and EF with both a provable lower
bound $\mathrm{LB}_R$ and a provable bridge $\mathrm B_R$.
Outcome: (A) is reduced to a strictly more local statement (Theorem~\ref{t:down}),
still OPEN. (B) fails in the predicted way: both routes for building
$R$-proofs from a SAT algorithm are blocked (Propositions~\ref{p:search},
\ref{p:sim}). Nothing here proves a lower bound for arbitrary algorithms.

\section{Target A: the bridge as an inductive invariant}
Encodings allow constants, so restriction $F\mapsto F|_{x=b}$ is a
polynomial-time map that does not increase encoding length. For a
polynomial-time $A$ and length $n$ define the propositional families
(each of size $\poly(n)$, with extension variables for the circuits of $A$
and of restriction):
\begin{align*}
\Down_{A,n}&:\ A(F)=0\ \rightarrow\ A(F|_{x=b})=0
 &&\text{(all $F$ of length $\le n$, all variables $x$, $b\in\{0,1\}$)},\\
\Base_{A,n}&:\ A(G)=0\ \rightarrow\ \operatorname{eval}(G)=0
 &&\text{(all closed $G$ of length $\le n$)}.
\end{align*}
Let $\mathrm{Rej}_A=\{F: A(F)=0\}$.

\begin{lemma}[Semantics; PROVED/EXACT]\label{l:sem}
$\Down_{A,n}\wedge\Base_{A,n}$ are tautologies iff $\mathrm{Rej}_A$ (at length
$\le n$) contains only UNSAT formulas.
\end{lemma}
\begin{proof}
($\Leftarrow$) UNSAT is closed under restriction and contains no true closed
formula. ($\Rightarrow$) If $F\in\mathrm{Rej}_A$ and $F(\alpha)=1$, restrict
along $\alpha$ one variable at a time; by $\Down$ every stage stays in
$\mathrm{Rej}_A$, and the closed end formula evaluates to $F(\alpha)=1$,
contradicting $\Base$.
\end{proof}

\begin{theorem}[Inductive-invariant form of $\mathrm B_{EF}$; PROVED/EXACT
modulo standard Cook translations]\label{t:down}
\[
\mathrm B_{EF}\iff\Bigl(\SAT\in\Pc\Rightarrow\exists A\ \text{poly-time with }
\mathrm{UNSAT}\subseteq\mathrm{Rej}_A,\ \text{and }\Down_{A,n},\Base_{A,n}
\text{ have }\poly(n)\text{-size EF proofs}\Bigr).
\]
\end{theorem}
\begin{proof}
($\Leftarrow$) Let $F$ be UNSAT, so $A(F)=0$, and the concrete evaluation
$A(\langle F\rangle)=0$ has a polynomial EF proof. Fix assignment variables
$\alpha$; set $G_0=F$, $G_k=G_{k-1}|_{x_k=\alpha_k}$ (circuits of size
$\poly(n)$). The $n$ instances of $\Down$ at $(G_{k-1},x_k,\alpha_k)$ chain to
$A(G_n)=0$; $\Base$ at $G_n$ gives $\operatorname{eval}(G_n)=0$; and
$\operatorname{eval}(G_n)=F(\alpha)$ is the translation of a
$\mathrm{PV}$-provable identity about restriction, hence has polynomial EF
proofs [Cook 1975]. Modus ponens yields a $\poly(n)$-size EF refutation of
$F$. (Each $\Down$ instance is the given proof with circuits substituted
for its free variables; EF is closed under such substitution up to
polynomial size.)
($\Rightarrow$) If $\SAT\in\Pc$ via a correct $A$ and $\mathrm B_{EF}$ holds,
EF is polynomially bounded; $\Down_{A,n},\Base_{A,n}$ are tautologies by
Lemma~\ref{l:sem}, so they have polynomial EF proofs.
\end{proof}

\paragraph{What this changes.} Compared with Theorem~4 of WR-BRIDGE-0001
(global correctness tautologies $\mathrm{Corr}_{A,n}$):
(i) Correctness on SAT inputs is no longer needed; only
$\mathrm{UNSAT}\subseteq\mathrm{Rej}_A$ (semantic completeness on the UNSAT
side) plus a \emph{one-step} invariance certificate.
(ii) The quantifier over assignments disappears from the certified
statement: $\Down$ relates $A$ on $F$ to $A$ on one restriction.
(iii) Exact structural correspondence with the user's DAG picture: the
restriction DAG of $F$ (nodes are partially reduced formulas, edges are
single substitutions) is the reduction DAG, and $\mathrm{Rej}_A$ must be a
\emph{forward-invariant set} of that DAG that excludes the ``unsafe'' true
closed formulas and contains every UNSAT start node. This is exactly the
shape of a barrier certificate in control theory (invariant, excludes
unsafe set, contains initial set) [Prajna--Jadbabaie 2004]. Classification:
EXACT OBJECT correspondence; no theorem transfers, because barrier
certificate theory supplies no bound on the \emph{proof size} of invariance.

\paragraph{Red team on Target A.}
\begin{enumerate}
\item \emph{Circularity.} $\neg\Down_{A,n}$ is an explicit UNSAT CNF
(after the Tseitin transformation). So Theorem~\ref{t:down} says EF is
polynomially bounded on all UNSAT formulas iff it is polynomially bounded on
the single family $\{\neg\Down_{A,n},\neg\Base_{A,n}\}$ for one suitable $A$.
This collapses the bridge to one family, but that family depends on the
unknown $A$. A polynomial SAT algorithm can \emph{decide} that family, but
deciding it does not give an EF proof. That gap is the bridge itself.
\item \emph{Trivial invariants are exponential.} The recursive
$A(F)=A(F|_{x=0})\vee A(F|_{x=1})$ makes $\Down$ EF-obvious but takes
exponential time. Polynomial-time algorithms that are known to be invariant in an
EF-provable way (unit propagation, 2-SAT implication graphs, Gaussian
elimination) all have $\mathrm{PV}$-provable correctness, so for them the
bridge is the known Cook result. No algorithm with polynomial time and
EF-short $\Down$ proofs is known on a domain that is hard for EF.
\item \emph{Status.} $\mathrm B_{EF}$: OPEN. Target A: reduced, not proved.
\end{enumerate}

\section{Target B: an intermediate system}
Candidates with unconditional lower bounds from the current literature:
Resolution (Haken), RAT$^-$, i.e.\ RAT without new variables (exponential
lower bound for the bit pigeonhole principle, Buss--Thapen 2019/2021),
bounded-depth and regular Res($\oplus$) (Efremenko--Garlik--Itsykson;
Alekseev--Itsykson; Byramji--Impagliazzo 2025 give
$\exp(\Omega(n^3/D^2))$ for depth-$D$ Res($\oplus$) proofs of BPHP;
Efremenko--Itsykson, Strong ETH for bounded-depth Res($\oplus$)).
For each of these, Proposition~2(c) of WR-BRIDGE-0001 gives
$\mathrm B_R\Leftrightarrow\Pc\ne\NP$. Given the lower-bound family
$H$, $\mathrm B_R$ is equivalent to the single-family statement
$\SAT\in\Pc\Rightarrow w_R(H_n)\le\poly(n)$, which must therefore be
derived by contradiction. There are two ways to obtain $R$-proofs from a
SAT algorithm $A$:

\begin{proposition}[Proof search is dead; PROVED/EXACT]\label{p:search}
Any procedure that uses $A$ only to \emph{find} $R$-proofs (for example,
self-reduction on ``an $R$-proof of size $\le s$ with prefix $p$ exists'')
outputs a proof of size $\le s$ only if one exists. On $H$ with
$w_R(H_n)>\poly(n)$ it therefore produces nothing short.
\end{proposition}

\begin{proposition}[Simulation requires extension; KNOWN THEOREM plus CANDIDATE]\label{p:sim}
The other route is Cook-style: simulate $A$'s computation inside $R$, as in
Theorem~\ref{t:down}. It needs lines that name the gates of $A$'s circuit,
i.e.\ new variables. Adding new variables to the candidate systems gives
systems p-equivalent to Extended Resolution: DRAT with new variables
p-simulates ER [Kiesl--Rebola-Pardo--Heule 2018], and ER $\equiv$ EF
[Cook--Reckhow 1979]. CANDIDATE claim (not proved): no system without
extension can carry out the simulation step for general polynomial-time
$A$. Feasible interpolation, which Resolution and Cutting Planes have,
is the natural tool for turning this into a theorem.
\end{proposition}
Conclusion: Target B fails as WR-BRIDGE-0001 predicted. Once a system can
run the bridge argument, it is EF-strength and has no known lower bound.

\section{Computational controls (COMPUTATIONAL)}
Script \path{experiments/claude_wr/certificate_closure.py}
(SHA-256 \texttt{45cd6a22087871449462c7b0b28d16c16f3c7085f183e8a9a61f6108890156c6},
seed 20261114, single stage, 4.0\,s, no chunking needed). Raw output:
\path{results/raw/claude_wr/2026-11-14_certificate_closure.json}.

\emph{2-SAT.} Random UNSAT 2-CNFs ($m=2n$). The decider's rejection
certificate (implication paths $x\Rightarrow\neg x\Rightarrow x$) was turned
into an explicit Resolution refutation, and every step was checked:
6, 11, 9, 16, 14 steps for $n=8,16,32,64,128$, always $\le 2m+1$. Resolution
contains this algorithm's certificates.

\emph{Tseitin, odd charge.} Gaussian elimination rejects with 3, 5, 5, 7
row XORs (each costs $O(m)$ bit operations). Exact minimum tree-like
Resolution leaves are 28, 76, 48, 128 for $K_4$, $K_{3,3}$, prism and $Q_3$
(6, 9, 9, 12 edge variables). Resolution does \emph{not} contain the affine
algorithm's certificates; asymptotically this is the Tseitin lower bound
[Urquhart 1987; Ben-Sasson--Wigderson 2001]. Small graphs are only
illustrations.

\paragraph{Certificate-closure ledger (KNOWN THEOREMS).}
For Boolean CSPs (Schaefer), the tractable classes Horn, dual-Horn and
bijunctive have rejection certificates that bounded-width Resolution can
check. The affine class needs parity reasoning: Res($\oplus$), or Polynomial
Calculus over $\mathbb F_2$. For general CSPs, bounded-width Resolution
refutes exactly the languages of bounded width [Barto--Kozik 2014]. A proof
system that satisfies the bridge must contain the certificates of
\emph{every} polynomial-time algorithm, including ones not yet known; this
is the precise sense in which $\mathrm B_R$ asks for a p-optimal-like $R$
[Krajíček--Pudlák 1989].

\section*{Literature check (this run)}
Res($\oplus$): ECCC TR24-128, TR25-039, TR25-118, TR25-188, TR26-007, CCC 2025;
Byramji--Impagliazzo 2025 (BPHP, depth-$D$); Efremenko--Itsykson, Strong ETH for
bounded-depth Res($\oplus$). DRAT/PR without new variables: Buss--Thapen,
arXiv:1909.00520, LMCS. No unrestricted Res($\oplus$) superpolynomial lower bound
and no EF lower bound found.

\section*{CORPUS REVIEWED}
WR-BRIDGE-0001 and its raw data; Codex 2026-11-11 SAT-polynomial audit
(bridge (12)); AGENTS.md; GROKPROOF 17:19 UTC note; repository search for
``inductive invariant'', ``Down'' and ``barrier certificate'' in agent reports
found no prior restriction-invariant form of the EF bridge.

\section*{DUPLICATION / CONTRADICTION AUDIT}
Theorem~\ref{t:down} refines Theorem~4 of WR-BRIDGE-0001 and is not
duplicated elsewhere in the corpus. Propositions~\ref{p:search} and \ref{p:sim}
are standard observations recorded for the directive. No contradictions.

\section*{WHAT CHANGED RELATIVE TO ROOT MAIN}
Root \texttt{main.tex} was not edited (this routine has no ownership region).

\section*{ADOPTED / REJECTED / RETIRED PRIOR COMPONENTS}
Retired: Target B as a route to an unconditional proof (any candidate $R$
with a known lower bound has $\mathrm B_R\Leftrightarrow\Pc\ne\NP$, and both
proof-generation routes are blocked). Adopted: Target A in the $\Down$ form.

\section*{STALE OR OVERSTATED CLAIMS FOUND}
None.

\section*{Next decisive target}
OPEN: assuming $\SAT\in\Pc$, construct a polynomial-time $A$ whose
rejection set contains UNSAT and whose one-step restriction invariance
$\Down_{A,n}$ has polynomial-size EF proofs. Alternatively, prove a
relativized or feasible-interpolation barrier showing that no
black-box construction from an arbitrary $A$ can do this.
\end{document}
% END_OF_VERIFIED_REPORT::CLAUDE::WR-BRIDGE-0002
